\documentclass{article}
\usepackage[latin1]{inputenc}
\usepackage{amssymb}
\usepackage{fancyvrb}
\usepackage{german}
\usepackage{a4wide}
\usepackage{epsfig}
\usepackage{fleqn}

\pagestyle{plain}

\def\pair(#1,#2){\langle #1, #2 \rangle}

\newcounter{aufgabe}

\renewcommand{\labelenumi}{(\alph{enumi})}

\newcommand{\exercise}{\vspace*{0.1cm}

\stepcounter{aufgabe}
\vspace*{0.3cm}

\noindent
\textbf{Aufgabe \arabic{aufgabe}}: }

\newcommand{\solution}{
\vspace*{0.3cm}

\noindent
\textbf{L\"osung}: }
\newcommand{\quoted}[1]{\;\mbox{``\texttt{#1}''}\;}


\begin{document}


\noindent
{\large  L\"osungen zu den Aufgaben zur Klausurvorbereitung}
\vspace{0.5cm}

\exercise
Die Grammatik $G = \langle \{ S \}, \{ \quoted{+}, \quoted{-}, \quoted{a} \}, R, S \rangle$
habe  die folgenden Regeln:
\[ S \rightarrow S\, S\, \quoted{+} \mid S\, S\, \quoted{-} \mid \quoted{a}. \]
\begin{enumerate}
\item Berechnen Sie die Mengen $\textsl{First}(S)$ und $\textsl{Follow}(S)$.
\item Berechnen Sie die Menge der SLR-Zust\"ande f\"ur diese Grammatik.
\item Berechnen Sie die Funktionen $\textsl{action}()$ und $\textsl{goto}()$ f\"ur diese Grammatik.
\item Berechnen Sie die Menge der LR-Zust\"ande f\"ur diese Grammatik.
\end{enumerate}

\solution
\begin{enumerate}
\item Es gilt offenbar
      \\[0.2cm]
      \hspace*{1.3cm}
      $\textsl{First}(S) = \{ a \}$ \quad und \quad
      $\textsl{Follow}(S) = \{ \quoted{+}, \quoted{-}, \quoted{a}, \symbol{36} \}$.
\item Wir erhalten die folgenden Zust\"ande:
      \begin{enumerate}
      \item Wir definieren $s_0 = \textsl{closure}\bigl(\{\widehat{S} \rightarrow \star S\}\bigr)$ und 
            finden
            \\[0.2cm]
            \hspace*{1.3cm}
            $s_0 = \{\; \widehat{S} \rightarrow \star\; S,\;
                      S \rightarrow \star \quoted{a},\; 
                      S \rightarrow \star \; S\; S \quoted{+},\; 
                      S \rightarrow \star \; S\; S \quoted{-}
            \}$.
      \item Wir definieren $s_1 = \textsl{goto}(s_0, S)$ und finden
            \\[0.2cm]
            \hspace*{1.3cm}
            $
            \begin{array}[t]{lcl}
            s_1 \;= & \{ & \widehat{S} \rightarrow S\; \star,             \\
                    &    & S \rightarrow \star \quoted{a},                \\
                    &    & S \rightarrow \star\; S\, S \quoted{+},        \\
                    &    & S \rightarrow \star\; S\, S \quoted{-},        \\
                    &    & S \rightarrow S \star S \quoted{+},            \\
                    &    & S \rightarrow S \star S \quoted{-}\;           \\
                    & \}.
            \end{array}
            $

      \item Wir definieren $s_2 = \textsl{goto}(s_1, S)$ und finden
            \\[0.2cm]
            \hspace*{1.3cm}
            $
            \begin{array}[t]{lcl}
            s_2 \;= & \{ & S \rightarrow \star \quoted{a},          \\
                    &    & S \rightarrow \star\; S\; S \quoted{+},  \\
                    &    & S \rightarrow \star\; S\; S \quoted{-},  \\
                    &    & S \rightarrow S \star S \quoted{+},    \\
                    &    & S \rightarrow S \star S \quoted{-},    \\
                    &    & S \rightarrow S\; S \star \quoted{+},  \\
                    &    & S \rightarrow S\; S \star \quoted{-}   \\
                    & \}.
            \end{array}
            $
      \item Wir definieren $s_3 = \textsl{goto}(s_2, \quoted{a})$ und finden
            \\[0.2cm]
            \hspace*{1.3cm}
            $s_3 = \{ S \rightarrow \quoted{a} \star \}$.
      \item Wir definieren $s_4 = \textsl{goto}(s_2, \quoted{+})$ und finden
            \\[0.2cm]
            \hspace*{1.3cm}
            $s_4 = \{ S \rightarrow S\, S \quoted{+} \star \}$.
      \item Wir definieren $s_5 = \textsl{goto}(s_2, \quoted{-})$ und finden
            \\[0.2cm]
            \hspace*{1.3cm}
            $s_5 = \{ S \rightarrow S \, S \quoted{-} \star \}$.
      \end{enumerate}
      \pagebreak

\item Damit erhalten wir f\"ur die Funktion $\textsl{action}()$ die folgende Tabelle:
      \begin{enumerate}
      \item $\textsl{action}(s_0, \quoted{a}) = \pair(\texttt{shift}, s_3)$
      \item $\textsl{action}(s_1, \quoted{\symbol{36}}) = \mathtt{accept}$
      \item $\textsl{action}(s_1, \quoted{a}) = \pair(\texttt{shift}, s_3)$
      \item $\textsl{action}(s_2, \quoted{+}) = \pair(\texttt{shift}, s_4)$
      \item $\textsl{action}(s_2, \quoted{-}) = \pair(\texttt{shift}, s_5)$
      \item $\textsl{action}(s_2, \quoted{a}) = \pair(\texttt{shift}, s_3)$
      \item $\textsl{action}(s_3, \quoted{\symbol{36}}) = \pair(\texttt{reduce}, S \rightarrow \quoted{a})$
      \item $\textsl{action}(s_3, \quoted{+}) = \pair(\texttt{reduce}, S \rightarrow \quoted{a})$
      \item $\textsl{action}(s_3, \quoted{-}) = \pair(\texttt{reduce}, S \rightarrow \quoted{a})$
      \item $\textsl{action}(s_3, \quoted{a}) = \pair(\texttt{reduce}, S \rightarrow \quoted{a})$
      \item $\textsl{action}(s_4, \quoted{\symbol{36}}) = \pair(\texttt{reduce}, S \rightarrow S\, S \quoted{+})$
      \item $\textsl{action}(s_4, \quoted{+}) = \pair(\texttt{reduce}, S \rightarrow S\, S \quoted{+})$
      \item $\textsl{action}(s_4, \quoted{-}) = \pair(\texttt{reduce}, S \rightarrow S\, S \quoted{+})$
      \item $\textsl{action}(s_4, \quoted{a}) = \pair(\texttt{reduce}, S \rightarrow S\, S \quoted{+})$
      \item $\textsl{action}(s_5, \quoted{\symbol{36}}) = \pair(\texttt{reduce}, S \rightarrow S\, S \quoted{-})$
      \item $\textsl{action}(s_5, \quoted{+}) = \pair(\texttt{reduce}, S \rightarrow S\, S \quoted{-})$
      \item $\textsl{action}(s_5, \quoted{-}) = \pair(\texttt{reduce}, S \rightarrow S\, S \quoted{-})$
      \item $\textsl{action}(s_5, \quoted{a}) = \pair(\texttt{reduce}, S \rightarrow S\, S \quoted{-})$
      \end{enumerate}
      F\"ur die Funktion $\textsl{goto}()$ finden wir:
      \begin{enumerate}
      \item $\textsl{goto}(s_0, S) = s_1$
      \item $\textsl{goto}(s_1, S) = s_2$
      \item $\textsl{goto}(s_2, S) = s_2$
      \end{enumerate}
\item Wir erhalten die folgenden Zust\"ande:
      \begin{enumerate}
      \item Wir setzen wieder 
            $s_0 = \textsl{closure}\Bigl(\bigl\{\widehat{S} \rightarrow \star S: \symbol{36}\bigr\}\Bigr)$ 
            und erhalten diesmal
            \\[0.2cm]
            \hspace*{1.3cm}
            $
            \begin{array}[t]{lcl}  
            s_0 \;= & \bigl\{ & \widehat{S} \rightarrow \star\; S: \quoted{\symbol{36}},                  \\
                    &         & S \rightarrow \star \quoted{a}: \{\quoted{\symbol{36}}, \quoted{a}\},       \\
                    &         & S \rightarrow \star\; S\; S \quoted{+}: \{\quoted{\symbol{36}}, \quoted{a}\}, \\
                    &         & S \rightarrow \star\; S\; S \quoted{-}: \{\quoted{\symbol{36}}, \quoted{a}\}  \\
                    & \bigr\}.
            \end{array}
            $
      \item Wir definieren $s_1 = \textsl{goto}(s_0, S)$ und erhalten
            \\[0.2cm]
            \hspace*{1.3cm}
            $
            \begin{array}[t]{lcl} 
         s_1 \; = & \bigl\{ & \widehat{S} \rightarrow S \star: \quoted{\symbol{36}},    \\
                  &         & S \rightarrow \star \quoted{a}: \{\quoted{+}, \quoted{-}, \quoted{a}\},  \\
                  &         & S \rightarrow \star\; S\; S \quoted{+}: \{\quoted{+}, \quoted{-}, \quoted{a}\}, \\
                  &         & S \rightarrow \star\; S\; S \quoted{-}: \{\quoted{+}, \quoted{-}, \quoted{a}\}, \\
                  &         & S \rightarrow S \star S \quoted{+}: \{\quoted{\symbol{36}}, \quoted{a}\},     \\
                  &         & S \rightarrow S \star S \quoted{-}: \{\quoted{\symbol{36}}, \quoted{a}\}      \\
                  & \bigr\}.
            \end{array}
            $
       \item Wir definieren $s_2 = \textsl{goto}(s_1, S)$ und erhalten
             \\[0.2cm]
             \hspace*{1.3cm}
             $
             \begin{array}[t]{lcl} 
          s_2 \;= & \bigl\{ & S \rightarrow \star \quoted{a}: \{\quoted{+}, \quoted{-}, \quoted{a}\},       \\
                  &         & S \rightarrow \star\; S\; S \quoted{+}: \{\quoted{+}, \quoted{-}, \quoted{a}\}, \\
                  &         & S \rightarrow \star\; S\; S \quoted{-}: \{\quoted{+}, \quoted{-}, \quoted{a}\}, \\
                  &         & S \rightarrow S \star S \quoted{+}: \{\quoted{+}, \quoted{-}, \quoted{a}\},   \\
                  &         & S \rightarrow S \star S \quoted{-}: \{\quoted{+}, \quoted{-}, \quoted{a}\},   \\
                  &         & S \rightarrow S\; S \star \quoted{+}: \{\quoted{\symbol{36}}, \quoted{a}\},   \\
                  &         & S \rightarrow S\; S \star \quoted{-}: \{\quoted{\symbol{36}}, \quoted{a}\}    \\
                  & \bigr\}.
             \end{array}
             $
       \item Wir definieren $s_3 = \textsl{goto}(s_2, S)$ und erhalten
             \\[0.2cm]
             \hspace*{1.3cm}
             $
             \begin{array}[t]{lcl} 
          s_3 \;= & \bigl\{ & S \rightarrow \star \quoted{a}: \{\quoted{+}, \quoted{-}, \quoted{a}\},       \\
                  &         & S \rightarrow \star\; S\; S \quoted{+}: \{\quoted{+}, \quoted{-}, \quoted{a}\}, \\
                  &         & S \rightarrow \star\; S\; S \quoted{-}: \{\quoted{+}, \quoted{-}, \quoted{a}\}, \\
                  &         & S \rightarrow S \star S \quoted{+}: \{\quoted{+}, \quoted{-}, \quoted{a}\},   \\
                  &         & S \rightarrow S \star S \quoted{-}: \{\quoted{+}, \quoted{-}, \quoted{a}\},   \\
                  &         & S \rightarrow S\; S \star \quoted{+}: \{\quoted{+}, \quoted{-}, \quoted{a}\}, \\
                  &         & S \rightarrow S\; S \star \quoted{-}: \{\quoted{+}, \quoted{-}, \quoted{a}\}  \\
                  & \bigr\}.
             \end{array}
             $
       \item Wir definieren $s_4 = \textsl{goto}(s_0, \quoted{a})$ und erhalten
             \\[0.2cm]
             \hspace*{1.3cm}
             $s_4 = \{ S \rightarrow \quoted{a} \star: \{\quoted{\symbol{36}}, \quoted{a}\} \}$.
       \item Wir definieren $s_5 = \textsl{goto}(s_2, \quoted{a})$ und erhalten 
             \\[0.2cm]
             \hspace*{1.3cm}
             $s_5 = \bigl\{ S \rightarrow \quoted{a} \star: \{\quoted{+}, \quoted{-}, \quoted{a}\} \bigr\}$
       \item Wir definieren $s_6 = \textsl{goto}(s_2, \quoted{+})$ und erhalten 
             \\[0.2cm]
             \hspace*{1.3cm}
             $s_6 = \bigl\{ S \rightarrow S\, S \quoted{+} \star: \{\quoted{\symbol{36}}, \quoted{a}\} \bigr\}$
       \item Wir definieren $s_7 = \textsl{goto}(s_3, \quoted{+})$ und erhalten 
             \\[0.2cm]
             \hspace*{1.3cm}
             $s_7 = \bigl\{ S \rightarrow S\, S \quoted{+} \star: \{\quoted{+}, \quoted{-}, \quoted{a}\} \bigr\}$
       \item Wir definieren $s_8 = \textsl{goto}(s_2, \quoted{-})$ und erhalten 
             \\[0.2cm]
             \hspace*{1.3cm}
             $s_8 = \bigl\{ S \rightarrow S\, S \quoted{-} \star: \{\quoted{\symbol{36}}, \quoted{a}\} \bigr\}$
       \item Wir definieren $s_9 = \textsl{goto}(s_3, \quoted{-})$ und erhalten 
             \\[0.2cm]
             \hspace*{1.3cm}
             $s_9 = \bigl\{ S \rightarrow S\, S \quoted{-} \star: \{\quoted{+}, \quoted{-}, \quoted{a}\} \bigr\}$
       \end{enumerate}
\end{enumerate}
\pagebreak

\exercise
Die Grammatik $G = \langle \{ A, B \}, \{ \quoted{u}, \quoted{x}, \quoted{y}, \quoted{z} \}, R, A \rangle$
habe  die folgenden Regeln:
\\[0.2cm]
\hspace*{1.3cm}
$\begin{array}[t]{lcl}
  A & \rightarrow & B \quoted{x} \\
    & \mid        & \quoted{y} B \quoted{z} \\
    & \mid        & \quoted{u} \quoted{z} \\
    & \mid        & \quoted{y} \quoted{u} \quoted{x} \\[0.1cm]
  B & \rightarrow &\quoted{u} 
 \end{array}
$
\\[0.2cm]
Bearbeiten Sie die folgenden Teilaufgaben:
\begin{enumerate}
\item \"Uberpr\"ufen Sie, ob die  diese Grammatik eine LL(1)-Grammatik ist und begr\"unden Sie
      Ihre Antwort.
\item \"Uberpr\"ufen Sie, ob die  diese Grammatik eine LL($*$)-Grammatik ist und begr\"unden Sie
      Ihre Antwort.
\item \"Uberpr\"ufen Sie, ob die  diese Grammatik eine SLR-Grammatik ist und begr\"unden Sie
      Ihre Antwort.
\end{enumerate}

\solution
\begin{enumerate}
\item Die Grammatik ist keine $LL(1)$-Grammatik, denn zwischen den beiden Regeln
      \\[0.2cm]
      \hspace*{1.3cm}
      $A \rightarrow \quoted{y} B \quoted{z}$ \quad und \quad
      $A \rightarrow \quoted{y} \quoted{u} \quoted{x}$ 
      \\[0.2cm]
      gibt es einen Konflikt, wir haben
      \\[0.2cm]
      \hspace*{1.3cm}
      $\textsl{First}(\quoted{y} B \quoted{z}) = \{ \quoted{y} \}$ \quad und \quad
      $\textsl{First}(\quoted{y} \quoted{u} \quoted{z}) = \{ \quoted{y} \}$ 
      \\[0.2cm]
      und damit folgt
      \\[0.2cm]
      \hspace*{1.3cm}
      $\textsl{First}(\quoted{y} B \quoted{z}) \cap \textsl{First}(\quoted{y} \quoted{u} \quoted{z}) = 
      \{ \quoted{y} \} \not= \{\}$.
\item Um zu \"uberpr\"ufen, ob die Grammatik eine LL($*$)-Grammatik ist, annotieren wir die Regeln
      wie im Skript beschrieben mit Zust\"anden und erhalten die folgenden annotierten Regeln:
      \begin{enumerate}
      \item $A \rightarrow _{\pair(0,1)} B _{\pair(1,1)} \quoted{x} _{\pair(2,1)}$
      \item $A \rightarrow _{\pair(0,2)} \quoted{y} _{\pair(1,2)} B _{\pair(2,2)} \quoted{z} _{\pair(3,2)}$
      \item $A \rightarrow _{\pair(0,3)} \quoted{u} _{\pair(1,3)} \quoted{z} _{\pair(2,3)}$
      \item $A \rightarrow _{\pair(0,4)} \quoted{y} _{\pair(1,4)} \quoted{u} _{\pair(2,4)} \quoted{x} _{\pair(3,4)}$
      \item $B \rightarrow _{\pair(5,0)} \quoted{u} _{\pair(6,0)}$ 
      \end{enumerate}
      Als Start-Zustand definieren wir den Zustand $\pair(0,0)$ und setzen
      $\textsl{start}(B) = \pair(4,0)$ sowie $\textsl{end}(B) = \pair(7,0)$.  Die Zustands-\"Ubergangs-Funktion des f\"ur
      einen LL($*)$-Parser zu konstruierenden nicht-deterministischen Automaten ist dann
      wie folgt gegeben:
      \begin{enumerate}
      \item Der Start-Zustand ist mit den Anfangs-Zust\"anden der einzelnen Regeln f\"ur $A$ verbunden,
            wir haben also
            \\[0.2cm]
            \hspace*{1.3cm} $\delta(\pair(0,0), \varepsilon) = \pair(0,1)$, \quad
                            $\delta(\pair(0,0), \varepsilon) = \pair(0,2)$, 
            \\[0.2cm]
            \hspace*{1.3cm} $\delta(\pair(0,0), \varepsilon) = \pair(0,3)$, \quad
                            $\delta(\pair(0,0), \varepsilon) = \pair(0,4)$.
      \item Aus der Regel $A \rightarrow B \quoted{x}$ erhalten wir die \"Uberg\"ange
            \\[0.2cm]
            \hspace*{1.3cm}
            $\delta(\pair(0,1), \varepsilon) = \pair(4,0)$, \quad
            $\delta(\pair(7,0), \varepsilon) = \pair(1,1)$, \quad
            $\delta(\pair(1,1), \quoted{x}) = \pair(2,1)$.
      \item Aus der Regel $A \rightarrow \quoted{y} B \quoted{z}$ erhalten wir die \"Uberg\"ange
            \\[0.2cm]
            \hspace*{1.3cm}
            $\delta(\pair(0,2), \quoted{y}) = \pair(1,2)$, \quad
            $\delta(\pair(1,2), \varepsilon) = \pair(4,0)$, \quad
            \\[0.2cm]
            \hspace*{1.3cm}
            $\delta(\pair(7,0), \varepsilon) = \pair(2,2)$, \quad
            $\delta(\pair(2,2), \quoted{z}) = \pair(3,2)$.
      \item Aus der Regel $A \rightarrow \quoted{u} \quoted{z}$ erhalten wir die \"Uberg\"ange
            \\[0.2cm]
            \hspace*{1.3cm}
            $\delta(\pair(0,3), \quoted{u}) = \pair(1,3)$, \quad
            $\delta(\pair(1,3), \quoted{z}) = \pair(2,3)$.
      \item Aus der Regel $A \rightarrow \quoted{y} \quoted{u} \quoted{x}$ erhalten wir die \"Uberg\"ange
            \\[0.2cm]
            \hspace*{1.3cm}
            $\delta(\pair(0,4), \quoted{y}) = \pair(1,4)$, \quad
            $\delta(\pair(1,4), \quoted{u}) = \pair(2,4)$, \quad
            $\delta(\pair(2,4), \quoted{x}) = \pair(3,4)$.
      \item Aus der Regel $B \rightarrow \quoted{u}$ erhalten wir die \"Uberg\"ange
            \\[0.2cm]
            \hspace*{1.3cm}
            $\delta(\pair(4,0), \varepsilon) = \pair(5,0)$, \quad
            $\delta(\pair(5,0), \quoted{u})  = \pair(6,0)$, \quad
            $\delta(\pair(6,0), \varepsilon) = \pair(7,0)$.
      \item Die Menge der akzeptierenden Zust\"ande ist durch die Endzust\"ande der Regeln f\"ur
            die Variable $A$ gegeben:
            \\[0.2cm]
            \hspace*{1.3cm}
            $\bigl\{ \pair(2,1), \pair(3,2), \pair(2,3), \pair(3,4) \bigr\}$
      \end{enumerate}
      Wir beginnen nun damit, diesen nicht-deterministischen Automaten in einen deterministischen 
      Automaten zu \"uberf\"uhren.  Der Start-Zustand des deterministischen  Automaten ist dann
      \\[0.2cm]
      \hspace*{1.3cm}
      $S_0 := 
       \bigl\{ \pair(0,0), \pair(0,1), \pair(0,2), \pair(0,3), \pair(0,4), \pair(4,0), \pair(5,0) \bigr\}$.
      \\[0.2cm]
      Wir definieren
      \\[0.2cm]
      \hspace*{1.3cm}
      $S_1 := \Delta(S_0, \quoted{u}) = 
      \bigl\{ \pair(6,0), \pair(7,0), \pair(1,1), \pair(2,2), \pair(1,3) \bigr\}$.
      \\[0.2cm]
      Weiter definieren wir
      \\[0.2cm]
      \hspace*{1.3cm}
      $S_2 := \Delta(S_1, \quoted{z}) = \bigl\{ \pair(3,2), \pair(2,3) \bigr\}$.
      \\[0.2cm]
      Diese Menge enth\"alt zwei akzeptierende Zust\"ande, ist aber nicht homogen, da diese
      Zust\"ande einen unterschiedlichen Index haben.  Folglich handelt es sich bei der obigen Grammatik
      nicht um eine LL($*$)-Grammatik.

      \textbf{Bemerkung}:
      \textsc{Antlr} kann aus der obigen Grammatik trotzdem einen Parser erzeugen, weil \textsc{Antlr}
      die Regel $B \rightarrow \quoted{u}$ expandiert und dann die folgende Grammatik analysiert:
      \\[0.2cm]
      \hspace*{1.3cm}
      $\begin{array}[t]{lcl}
        A & \rightarrow & \quoted{u} \quoted{x} \\
          & \mid        & \quoted{y} \quoted{u} \quoted{z} \\
          & \mid        & \quoted{u} \quoted{z} \\
          & \mid        & \quoted{y} \quoted{u} \quoted{x} 
        \end{array}
       $
       \\[0.2cm]
       Die so umgeformte Grammatik ist eine LL($*$)-Grammatik.
\item Die angegebene Grammatik ist keine SLR-Grammatik.  Um das zu sehen, erweitern wir die Grammatik
      um die Regel $\widehat{S} \rightarrow A$ und berechnen den Zustand
      \\[0.2cm]
      \hspace*{1.3cm}
      $s_0 = \textsl{closure}( \{ \widehat{S} \rightarrow A \})$.
      \\[0.2cm]
      Wir finden
      \\[0.2cm]
      \hspace*{1.3cm}
      $
      \begin{array}[t]{lcl}  
       s_0 \;= & \{ & S \rightarrow \star\; A,                              \\
               &    & A \rightarrow \star\; B \quoted{x},                   \\
               &    & A \rightarrow \star \quoted{y} B \quoted{z},          \\
               &    & A \rightarrow \star \quoted{u} \quoted{z},            \\
               &    & A \rightarrow \star \quoted{y} \quoted{u} \quoted{x}, \\
               &    & B \rightarrow \star \quoted{u}                        \\
               & \}. 
      \end{array}
      $
      \\[0.2cm]
      Wir berechnen nun $\textsl{goto}(s_0, \quoted{u})$ und erhalten
      \\[0.2cm]
      \hspace*{1.3cm}
      $s_1 = \{ A \rightarrow \quoted{u} \star \quoted{z},\;
                B \rightarrow \quoted{u} \star
             \}
      $.
      \\[0.2cm]
      Bei der Berechnung von $\textsl{action}(s_1, \quoted{z})$ tritt nun eine Shift-Reduce-Konflikt auf,
      denn es gilt
      \\[0.2cm]
      \hspace*{1.3cm}
      $\textsl{follow}(B) = \{ \quoted{x}, \quoted{z} \}$.

      \noindent
      \textbf{Bemerkung}:  Die in der Aufgabe angegebene Grammatik ist sowohl eine LR-Grammatik als
      auch eine LALR-Grammatik.  Letzteres l\"asst sich mit \textsl{Bison} oder \textsl{JavaCup} nachweisen 
      und Ersteres folgt aus der Tatsache, dass jede LALR-Grammatik auch eine LR-Grammatik ist.
\end{enumerate}
\pagebreak

\exercise
Wir definieren \emph{geschachtelte Listen} rekursiv als solche Listen, 
deren Elemente nat\"urliche Zahlen oder geschachtelte Listen sind.  Die Elemente in
geschachtelten Listen sollen durch Kommata getrennt werden und die Listen selber sollen
durch die eckigen Klammern ``\texttt{[}'' und ``\texttt{]}'' begrenzt sein.
Beispiele f\"ur geschachtelte Listen sind also:
\begin{enumerate}
\item \texttt{[ 1, [ 1, [], [ 2, 3]], 7]}
\item \texttt{[ [], [ [], [], [ 4], 5], []]}
\end{enumerate}
L\"osen Sie die folgenden Teilaufgaben:
\begin{enumerate}
\item Geben Sie eine Grammatik f\"ur geschachtelte Listen an.
\item Geben Sie die EP-Spezifikation von Klassen an, mit deren Hilfe sich
      geschachtelte Listen repr\"asentieren lassen.
\item Geben Sie einen \textsl{JavaCup}-Parser an, der eine geschachtelte Liste einliest
      und einen abstrakten Syntax-Baum der Liste berechnet.

      \textbf{Hinweis:} In der Klausur k\"onnen Sie sp\"ater davon ausgehen, dass ein geeigneter
      \textsl{JFlex}-Scanner bereits gegeben ist, aber bei dieser Aufgabe sollen Sie den
      Scanner ebenfalls erstellen, damit Sie Ihre L\"osung auch testen k\"onnen.
\end{enumerate}

\solution
\begin{enumerate}
\item Die Grammatik 
      $G = \langle \{S,L,N,E\}, \{ \quoted{[},\quoted{]}, \quoted{,}\textsc{Number} \}, R, S \rangle$, 
      deren Regeln durch
      \\[0.2cm]
      \hspace*{1.3cm}
      $
      \begin{array}[t]{lcl}
        S & \rightarrow & \quoted{[} L \quoted{]}         \\[0.2cm]
        L & \rightarrow & N                               \\
          & \mid        & \varepsilon                     \\[0.2cm]
        N & \rightarrow & N \quoted{,} E                  \\
          & \mid        & E                               \\[0.2cm]
        E & \rightarrow & S                               \\
          & \mid        & \textsc{Number}                 \\
      \end{array}
      $
      \\[0.2cm]
      gegeben sind, leistet das Gew\"unschte.  Hier steht $L$ f\"ur eine ungeklammerte Liste, 
      deren Elemente durch Kommata
      getrennt sind,  $N$ steht f\"ur eine ungeklammerte nicht-leere Liste, deren Elemente durch Kommata getrennt
      sind und $E$ steht f\"ur ein Listen-Element, ist also entweder eine in eckigen Klammern
      eingefasste Liste oder eine Zahl.
\item Wir definieren eine abstrakte Klasse \texttt{Element}, die sowohl Listen als auch Zahlen
      umfasst:
      \begin{Verbatim}[ frame         = lines, 
                        framesep      = 0.3cm, 
                        labelposition = bottomline,
                        numbers       = left,
                        numbersep     = -0.2cm,
                        xleftmargin   = 0.8cm,
                        xrightmargin  = 0.8cm,
                      ]
      Element = MyList(List<Element> listExpr) 
              + MyNumber(Integer number);
      \end{Verbatim}
      \pagebreak

\item Die \textsl{JavaCup}-Spezifikation des Parsers hat die folgende Form:
      \begin{Verbatim}[ frame         = lines, 
                        framesep      = 0.3cm, 
                        labelposition = bottomline,
                        numbers       = left,
                        numbersep     = -0.2cm,
                        xleftmargin   = -0.3cm,
                        xrightmargin  = 0.0cm,
                      ]
    import java_cup.runtime.*;
    import java.util.*;
    
    terminal         OPEN, CLOSE, COMMA;
    terminal Integer NUMBER;
    
    nonterminal Element       s;
    nonterminal List<Element> l, n;
    nonterminal Element       e;
    
    s ::= OPEN l:e CLOSE  
          {: List<Listen> list = new LinkedList<Listen>();
             RESULT = new MyList(e);
          :};
    
    l ::= n:l                {: RESULT = l; :}
       |  /* epsilon */      {: List<Element> l = new LinkedList<Element>(); 
                                RESULT = l; 
                             :}
       ;
    
    n ::= n:l COMMA e:x      {: l.add(x); RESULT = l; :}
       |  e:x                {: List<Element> l = new LinkedList<Element>(); 
                                l.add(x);
                                RESULT = l; 
                             :}
       ;
    
    e ::= s:x                {: RESULT = x;               :}
       |  NUMBER:x           {: RESULT = new MyNumber(x); :}
       ;
      \end{Verbatim}
      \pagebreak

      Der dazugeh\"orige Scanner ist wie folgt definiert:
      \begin{Verbatim}[ frame         = lines, 
                        framesep      = 0.3cm, 
                        labelposition = bottomline,
                        numbers       = left,
                        numbersep     = -0.2cm,
                        xleftmargin   = -0.3cm,
                        xrightmargin  = 0.0cm,
                ]
    import java_cup.runtime.*;
          
    %%
       
    %char
    %line
    %column
    %cup
       
    %{   
        private Symbol symbol(int type) {
            return new Symbol(type, yychar, yychar + yylength());
        }
        
        private Symbol symbol(int type, Object value) {
            return new Symbol(type, yychar, yychar + yylength(), value);
        }
    %}
       
    %%
       
    "["      { return symbol( sym.OPEN  ); }
    "]"      { return symbol( sym.CLOSE ); }
    ","      { return symbol( sym.COMMA ); }
    
    [0-9]+   { return symbol( sym.NUMBER, new Integer(yytext()) ); }
    
    [ \t\n]  { /* skip white space */ }   
    
    [^]      { throw new Error("Illegal character"); }
       \end{Verbatim}

\end{enumerate}
\pagebreak


\exercise
Der Typ $\texttt{list}(T)$ sei wie folgt definiert:
\\[0.2cm]
\hspace*{1.3cm}
\texttt{type list(X) := nil + cons(X, list(X));}
\\[0.2cm]
Die Funktion \texttt{addLast} habe die folgende Signatur:
\\[0.2cm]
\hspace*{1.3cm}
\texttt{signature addLast: list(T) * T -> list(T);}
\\[0.2cm]
und die Variablen \texttt{x} und \texttt{z} haben den Typ \texttt{int}.
\begin{enumerate}
\item Berechnen Sie 
      \\[0.2cm]
      \hspace*{1.3cm}
      $\textsl{typeEqs}\bigl(\texttt{addLast(cons(x, nil), z): list(int)}\bigr)$.
\item L\"osen Sie die in Teil (a) berechneten Typ-Gleichungen.
\end{enumerate}

\solution
\begin{enumerate}
\item Wir berechnen zun\"achst die Typ-Gleichungen nach der im Skript angegebenen Definition.
      \[
      \begin{array}[t]{cl}
           & \textsl{typeEqs}(\texttt{addLast(cons(x, nil), z): list(int)}) \\[0.2cm]
         = &       \{ \mathtt{list}(T) = \mathtt{list(int)} \} 
              \cup \textsl{typeEqs}(cons(x, nil): \mathtt{list}(T))
              \cup \textsl{typeEqs}(z: T)                     \\[0.2cm]
         = &       \{ \mathtt{list}(T) = \mathtt{list(int)} \} 
              \cup \{ \texttt{list}(S) = \mathtt{list}(T) \}\; \cup \\
           &       \textsl{typeEqs}(x: S)
              \cup \textsl{typeEqs}(\mathtt{nil}: \mathtt{list}(S))
              \cup \textsl{typeEqs}(z: T)                  \\[0.2cm]
         = &       \{ \mathtt{list}(T) = \mathtt{list(int)},\; 
                      \texttt{list}(S) = \mathtt{list}(T),\;           \mathtt{int} = S,\;
                      \mathtt{list}(R) = \mathtt{list}(S),\;
                      \mathtt{int} = T \}

      \end{array}
      \]
\item Wir l\"osen die oben berechneten Typ-Gleichungen nach dem im Skript angegebenen Verfahren.
      \\[0.2cm]
      \hspace*{-0.3cm}  
      $\begin{array}[t]{cl}
           &  \bigl\langle \{  \mathtt{list}(T) = \mathtt{list(int)},\; 
                      \texttt{list}(S) = \mathtt{list}(T),\;           \mathtt{int} = S,\;
                      \mathtt{list}(R) = \mathtt{list}(S),\;
                      \mathtt{int} = T \},\; [] \bigr\rangle \\[0.2cm]
         \leadsto &  \bigl\langle \{ T = \mathtt{int},\; 
                      \texttt{list}(S) = \mathtt{list}(T),\;           \mathtt{int} = S,\;
                      \mathtt{list}(R) = \mathtt{list}(S),\;
                      \mathtt{int} = T \},\; [] \bigr\rangle \\[0.2cm]
         \leadsto &  \bigl\langle \{ 
                      \texttt{list}(S) = \mathtt{list}(\mathtt{int}),\;           
                      \mathtt{int} = S,\;
                      \mathtt{list}(R) = \mathtt{list}(S),\;
                      \mathtt{int} = T \},\; [ T \mapsto \mathtt{int}] \bigr\rangle \\[0.2cm]
         \leadsto &  \bigl\langle \{ 
                      S = \mathtt{int},\;           
                      \mathtt{int} = S,\;
                      \mathtt{list}(R) = \mathtt{list}(S),\;
                      \mathtt{int} = \mathtt{int} \},\; [ T \mapsto \mathtt{int}] \bigr\rangle \\[0.2cm]
         \leadsto &  \bigl\langle \{            
                      \mathtt{int} = \mathtt{int},\;
                      \mathtt{list}(R) = \mathtt{list}(\mathtt{int}),\;
                      \mathtt{int} = \mathtt{int} \},\; 
                      [ T \mapsto \mathtt{int},\; S \mapsto \mathtt{int}] \bigr\rangle \\[0.2cm]
         \leadsto &  \bigl\langle \{            
                      \mathtt{list}(R) = \mathtt{list}(\mathtt{int}),\;
                      \mathtt{int} = \mathtt{int} \},\; 
                      [ T \mapsto \mathtt{int},\; S \mapsto \mathtt{int}] \bigr\rangle \\[0.2cm]
         \leadsto &  \bigl\langle \{            
                      R = \mathtt{int},\;
                      \mathtt{int} = \mathtt{int} \},\; 
                      [ T \mapsto \mathtt{int},\; S \mapsto \mathtt{int}] \bigr\rangle \\[0.2cm]
         \leadsto &  \bigl\langle \{            
                      \mathtt{int} = \mathtt{int} \},\; 
                      [ T \mapsto \mathtt{int},\; S \mapsto \mathtt{int},\; R \mapsto \mathtt{int}] \bigr\rangle \\[0.2cm]
         \leadsto &  \bigl\langle \{ \},\; 
                      [ T \mapsto \mathtt{int},\; S \mapsto \mathtt{int},\; R \mapsto \mathtt{int}] \bigr\rangle
         \end{array}
      $
      \\[0.2cm]
      Damit ist die Substitution 
      $[ T \mapsto \mathtt{int},\; S \mapsto \mathtt{int},\; R \mapsto \mathtt{int}]$ eine 
      L\"osung der Typ-Gleichungen und wir k\"onnen folgern, dass der Term tats\"achlich den angegebenen Typ
      hat.
\end{enumerate}
\pagebreak

\exercise
Nehmen Sie an, dass die im Skript eingef\"uhrte Sprache \textsl{Integer}-\texttt{C} um eine 
\texttt{do-while}-Schleife erweitert werden soll, deren Syntax durch die folgende Grammatik-Regel gegeben ist:
\\[0.2cm]
\hspace*{1.3cm}
$\textsl{statement} \rightarrow \quoted{do} \textsl{statement} \quoted{while}
 \quoted{(} \textsl{boolExpr} \quoted{)}$.
\\[0.2cm]
Die Semantik dieses Konstruktes soll mit der Semantik des entsprechenden Konstruktes in
der Sprache \texttt{C} \"ubereinstimmen.
\begin{enumerate}
\item Geben Sie eine Gleichung an, die beschreibt, wie eine \texttt{do-while}-Schleife in
      \textsl{Java-Byte-Code} \"ubersetzt werden kann.
\item Geben Sie die Methode $\textsl{compile}()$ an, die das entsprechende
      Konstrukt \"ubersetzt.  Gehen Sie dabei davon aus, dass Sie diese Methode innerhalb
      einer Klasse \texttt{DoWhile} implementieren, wobei diese Klasse f\"ur \textsc{Ep}
      wie folgt spezifiziert ist:
      \\[0.2cm]
      \hspace*{1.3cm}
      $\texttt{Statement} \;=\; \cdots \;+\; \texttt{DoWhile}(\texttt{Statement}\; stmnt,\;
      \mathtt{BoolExpr}\; cond)\; +\; \cdots; $      
\end{enumerate}

\solution
\begin{enumerate}
\item Die \"Ubersetzung einer \texttt{do-while}-Schleife der Form
      \[ \texttt{while}\;(\textsl{cond})\;\textsl{statement} \]
      orientiert sich an der folgenden Spezifikation:
      \[
      \begin{array}[t]{lcl}
        \textsl{compile}\bigl(\mathtt{do}\; \textsl{stmnt}\; \texttt{while}\; (\textsl{cond})\;\bigr) 
        & = & 
        [\;\textsl{loop}\texttt{:}\;]            \\
        & + & \textsl{stmnt}.\textsl{compile}()    \\    
        & + & \textsl{cond}.\textsl{compile}()         \\
        & + & [\;\texttt{ifeq}\;\textsl{next}\;] \\
        & + & [\;\texttt{goto}\;\textsl{loop}\;]    \\    
        & + & [\;\textsl{next}\texttt{:}\;]         
      \end{array}
      \]
\item Die Methode $\textsl{compile}()$ kann in der Klasse \texttt{DoWhile} wie folgt implementiert werden:
      \begin{Verbatim}[ frame         = lines, 
                        framesep      = 0.3cm, 
                        labelposition = bottomline,
                        numbers       = left,
                        numbersep     = -0.2cm,
                        xleftmargin   = 0.0cm,
                        xrightmargin  = 1.1cm,
                      ]
      public List<AssemblerCmd> compile() {
          List<AssemblerCmd> result = new LinkedList<AssemblerCmd>();
          LABEL        loopLabel = new LABEL();
          LABEL        nextLabel = new LABEL();
          AssemblerCmd ifeq      = new IFEQ(nextLabel.getLabel());
          AssemblerCmd gotoLoop  = new GOTO(loopLabel.getLabel());
          result.add(loopLabel);
          result.addAll(mStmnt.compile());
          result.addAll(mCond.compile());
          result.add(ifeq);
          result.add(gotoLoop);
          result.add(nextLabel);
          return result;
        }
      \end{Verbatim}
\end{enumerate}

\pagebreak

\noindent
\framebox{\epsfig{file=danger.eps,scale=0.05}}
\exercise
Zeigen Sie, dass die wie folgt definierte Menge $M$ 
\\[0.2cm]
\hspace*{1.3cm}
$M := \{ x \in \mathbb{R} \mid 0 \leq x \leq 1 \}$
\\[0.2cm] 
nicht abz\"ahlbar ist.

\noindent
\textbf{Hinweis}: Nehmen Sie an, dass Sie eine Funktion
\\[0.2cm]
\hspace*{1.3cm}
$f: \mathbb{N} \rightarrow M$
\\[0.2cm]
gibt, so dass es f\"ur jedes $x \in M$ eine Zahl $n \in \mathbb{N}$ gibt, f\"ur die
$x = f(n)$ ist.  F\"uhren Sie diese Annahme zum Widerspruch, indem Sie eine Zahl 
$\gamma \in M$ konstruieren, f\"ur die 
\\[0.2cm]
\hspace*{1.3cm}
$f(n) \not= \gamma$ \quad f\"ur alle $n \in \mathbb{N}$
\\[0.2cm]
gilt.  Zur Konstruktion von $\gamma$ ist es sinnvoll die Hilfsfunktion
\\[0.2cm]
\hspace*{1.3cm}
$\textsl{digit}: M \times \mathbb{N} \rightarrow \{ 0, 1, 2, 3, 4, 5, 6,7, 8, 9\}$
\\[0.2cm]
zu benutzen, die so definiert ist, dass
$\textsl{digit}(x,k)$
die $k$-te Ziffer der Zahl $x$ ist.  Beispielsweise gilt
\\[0.2cm]
\hspace*{1.3cm}
$\textsl{digit}(0.125, 1) = 1$, \quad
$\textsl{digit}(0.125, 2) = 2$, \quad
$\textsl{digit}(0.125, 3) = 5$, \quad und 
\\[0.2cm]
\hspace*{1.3cm}
$\textsl{digit}(0.125, k) = 0$, \quad f\"ur alle $k > 3$.

\solution
Wir definieren die Zahl $\gamma$ indem wir die einzelnen Ziffern von $\gamma$ durch
\\[0.2cm]
\hspace*{1.3cm}
$\textsl{digit}(\gamma,n) = \bigl(\textsl{digit}\bigl(f(n),n\bigr) + 1\bigr) \;\mathtt{mod}\; 10$
\\[0.2cm]
definieren.  Den Widerspruch erhalten wir nun aus der Tasache, dass es eine Zahl $n_0 \in \mathbb{N}$ geben 
muss, f\"ur die
\\[0.2cm]
\hspace*{1.3cm}
$f(n_0) = \gamma$
\\[0.2cm]
gilt.  Betrachten wir nun die $n_0$-te Ziffer von $\gamma$, so sehen wir
\\[0.2cm]
\hspace*{1.3cm}
$
\begin{array}[t]{lcl}
\textsl{digit}(f(n_0), n_0) & = & \textsl{digit}(\gamma, n_0)                          \\[0.2cm]
                            & = & (\textsl{digit}(f(n_0),n_0) + 1) \;\mathtt{mod}\; 10. 
\end{array}
$
\\[0.2cm]
Da f\"ur jede Ziffer $x \not= (x + 1) \;\mathtt{mod}\; 10$ gilt, ist dies ein Widerspruch, der uns zeigt,
dass die Menge $M$ nicht abz\"ahlbar ist.
\vspace{0.3cm}

\noindent
\framebox{\epsfig{file=danger.eps,scale=0.05}}
\exercise
Aus wieviel Zeichen muss der kleinste Ausdruck $e$ mindestens bestehen, f\"ur den
\\[0.2cm]
\hspace*{1.3cm}
$\textsl{ershov}(e) = 30$
\\[0.2cm]
gilt?

\solution
Wir bezeichnen die Anzahl der Zeichen, die der kleineste Ausdruck mit Ershov-Zahl $n$ haben muss,
mit $a_n$.  Dann haben wir $a_1 = 1$ und 
\\[0.2cm]
\hspace*{1.3cm}
$a_{n+1} = 2 * a_n + 1$
\\[0.2cm]
Diese Rekurrenz-Gleichung hat die L\"osung $a_{n} = 2^n - 1$, was wir leicht durch Induktion nachweisen k\"onnten.
Setzen wir f\"ur $n$ den Wert $30$ ein, so sehen wir, dass der kleinste Ausdruck aus
\\[0.2cm]
\hspace*{1.3cm}
$2^{30} - 1 = 1\,073\,741\,823$
\\[0.2cm]
Zeichen besteht.  Dieser Ausdruck h\"atte daher eine Gr\"o{\ss}e von knapp einem Gigabyte und wir k\"onnen
schlie{\ss}en, dass in der Praxis 30 Register f\"ur die Auswertung von Ausdr\"ucken ausreichen.

\end{document}

%%% Local Variables: 
%%% mode: latex
%%% TeX-master: t
%%% End: 
